\chapter{Cosets, stress splitting, and the near-two frontier}
\label{ch:cosets}

A central-charge-eight construction naturally invites a smaller theory:
can part of its stress tensor be removed while preserving Haagerup
symmetry? This chapter explains both a successful charge-six reduction
and the obstacles in the charge-two direction. It also constructs an
irrational stress of charge $2.025004294\ldots$ and follows the precise
distance between that chiral construction and a local full CFT.

The distinction is structural. Finding a Virasoro vector is an OPE
problem inside one chiral algebra. A full reduced theory additionally
needs a state space, a modular pairing, convergent positive OPEs, and
coherent defect junctions. None of those later requirements can be
inferred from a central charge being close to two.

\section{A commuting stress is not automatically a tensor factor}
\index{coset!commuting stress}

Let $V$ be unitary, with stress $\omega$ and translation operator
$D=L^V_{-1}$. Write $L^e_n=e_{(n+1)}$ for the modes of a
PCT-real weight-two vector $e$. Here PCT-real means fixed by the
antilinear involution of Chapter~\ref{ch:vertex_algebras}, with the
positive form transported in the chosen coordinate gauge.
\begin{definition}[Semiconformal Virasoro subalgebra]
\index{semiconformal vector}
The vector $e$ is semiconformal if the vertex subalgebra
$\Vir(e)\subset V$ has $e$ as its conformal vector and
\[
 L^e_n\big|_{\Vir(e)}=L^V_n\big|_{\Vir(e)}\quad(n\geq0),
 \qquad L^e_{-1}\big|_{\Vir(e)}=D\big|_{\Vir(e)}.
\]
These are equalities on the generated subalgebra; they are not
equalities on all of $V$. The last equality specifies compatibility
with the inherited translation operator.
\end{definition}

This definition explains why $f=\omega-e$ is a commuting stress.
For every $a\in\Vir(e)$ and $m\geq0$, the restrictions give
$f_{(m)}a=(L^V_{m-1}-L^e_{m-1})a=0$. Skew symmetry then gives
$a_{(m)}f=0$ as well. Thus $e$ and $f$ have no singular mixed OPE.
Subtracting their contributions from the ambient Virasoro OPE shows
that $f$ is a Virasoro vector of charge $c_V-c_e$; its translation
on $\Vir(f)$ is again the inherited $D$. In particular,
\[
 e(z)e(w)\sim\frac{c_e/2}{(z-w)^4}
       +\frac{2e(w)}{(z-w)^2}+\frac{\partial e(w)}{z-w},
 \qquad e(z)f(w)\text{ is regular},
\]
and $c_f=c_V-c_e$.

The commutant
\[
 \Com(A,V)=\{v\in V:a_{(n)}v=0
                       \text{ for all }a\in A,\ n\geq0\}
\]
collects fields whose OPEs with $A$ are regular. In the unitary
setting it is often an energy-zero space for a complementary stress.
This gives an actual chiral subalgebra with a definite positive form.
It does not imply that $V$ is the tensor product of the two
commutants. Nonvacuum paired modules may be needed to reconstruct $V$.

\begin{example}[A lattice extension with two commuting stresses]
The plane lattice $L_0$ of Chapter~\ref{ch:factors} and its orthogonal
rank-six lattice $N$ give commuting stress tensors inside $V_{E_8}$.
Their central charges add to eight. Nevertheless the $E_8$ lattice
contains glue cosets beyond $L_0\oplus N$, so
$V_{E_8}$ is a finite extension of $V_{L_0}\otimes V_N$, not that
tensor product itself. The extra paired sectors are exactly what a
naive stress subtraction would overlook.
\end{example}

There is also a categorical obstruction to an actual rational
$c=2$ times $c=6$ factorization of $U$. Its simple quantum dimensions
are $1,d^2$, four copies of $s=1+d^2$, and six copies of $3d$.
The unit is its only invertible simple, and
\[
 (3d)^2>s.
\]
If $\Rep(U)=\mathcal D\boxtimes\mathcal E$ with both factors
nontrivial, neither could be pointed. Choose nonunit simples $X,Y$
from the two factors. Each has dimension at least $3d$, whereas the
simple $X\boxtimes Y$ would exceed the largest allowed dimension $s$.
Thus no such categorical factorization exists. If one VOA tensor
factor were holomorphic, its central charge would be a multiple of
eight, incompatible with charges two or six. This excludes a tensor
product, while allowing proper conformal extensions and commutants.

\section{The concrete charge-six commutant}
\index{charge-six construction}

Recall
\[
 L_0=\Z\ell+\Z v,
 \qquad G_{L_0}=\begin{pmatrix}2&3\\3&24\end{pmatrix},
 \qquad\det L_0=39.
\]
Let $P$ be its real plane in the $E_8$ lattice $\Lambda$. Since
$39$ is squarefree, the displayed lattice is saturated:
$L_0=\Lambda\cap P$. Indeed an integral overlattice of index $r$
would have determinant $39/r^2$, forcing $r=1$.

Let $N=\Lambda\cap P^\perp$. Its rank is six and its determinant
is also $39$. An explicit description is
\begin{equation}\label{eq:c-rank6-lattice}
 N=(A_2\oplus A_2\oplus A_1\oplus\langle78\rangle)
        +\Z(\omega_1,\omega_2,\alpha/2,\gamma_{78}/6),
\end{equation}
Here $\omega_i$ is a fundamental weight of the $i$th $A_2$ factor,
$\alpha$ generates $A_1$, and $\gamma_{78}$ generates the rank-one
lattice $\langle78\rangle$. Their squared norms are
$\omega_i^2=2/3$, $\alpha^2=2$, and $\gamma_{78}^2=78$.
The glue vector has squared norm
\[
 \frac23+\frac23+\frac12+\frac{78}{36}=4.
\]
It has order six modulo the direct-sum lattice. Its index-six
extension has determinant
$3\cdot3\cdot2\cdot78/6^2=39$, as required.

Define
\begin{equation}\label{eq:c-C6}
 C_6=\Com(V_{L_0},U).
\end{equation}
Positive lattice energy gives an explicit ambient realization,
\[
 C_6=U\cap V_N.
\]
To see it, a state commuting with the plane Heisenberg algebra has
neither plane oscillator excitations nor plane momentum. Such states
are exactly $V_N$. Conversely, commuting with the full plane lattice
is no stronger on this subspace. Its stress is the perpendicular
free-boson stress, of charge six.

The pair is mutual:
\[
 \Com(C_6,U)=V_{L_0}.
\]
The key is that $C_6$ contains its full perpendicular stress. A vector
commuting with that stress has zero perpendicular energy, so has no
perpendicular momenta or oscillators. Saturation identifies the
remaining lattice states with precisely $V_{L_0}$. This is an
all-grade argument, not a low-weight count.

The regular lattice-coset theorem now applies: a strongly rational
VOA with this positive-definite lattice double commutant has a
strongly rational Heisenberg/lattice coset. Thus $C_6$ is strongly
rational. This uses a specific theorem with verified lattice
hypotheses; rationality of arbitrary cosets is not being assumed.

\section{Why the Haagerup extension data survive at charge six}

Let $P_{39}=\Rep(V_{L_0})$. It is the nondegenerate pointed category
on $\Z/39$ with quadratic form $q(n)=n^2/39$. The actual core
$A_L$ contains every one of these lattice labels. The mirror theorem
for the rational mutual commutants therefore identifies their branching
subcategory inside $\Rep(C_6)$ with $P_{39}^{\rm rev}$.
The superscript $\mathrm{rev}$ reverses the braiding and twist;
reversing tensor order is instead denoted $\mathrm{op}$ below.
Nondegeneracy splits off this pointed factor:
\[
 \Rep(C_6)=P_{39}^{\rm rev}\boxtimes\mathcal D.
\]
The extension $U$ of $V_{L_0}\otimes C_6$ is the diagonal algebra
on the two pointed factors, tensored with the unit in $\mathcal D$.
Condensing that diagonal algebra leaves $\mathcal D$, while the
extension correspondence says that it leaves $\Rep(U)$. Hence
\begin{equation}\label{eq:c-C6-category}
 \Rep(C_6)\simeq P_{39}^{\rm rev}\boxtimes\Rep(U).
\end{equation}
In the Haagerup construction its rank is $39\cdot12=468$.
The reverse pointed factor has chiral central charge six modulo eight,
consistent with the actual charge-six stress.

Let $A$ denote the actual algebra of the original $E_8/U$ extension.
Transitivity of extension algebras identifies the $E_8$ algebra over
$V_{L_0}\otimes C_6$ as
\[
 \operatorname{Diag}(P_{39})\boxtimes A.
\]
Taking the $V_{L_0}$ vacuum multiplicity leaves $\one\boxtimes A$.
The same vacuum multiplicity in $V_{E_8}$ is $V_N$, so the actual
extension $V_N/C_6$ realizes this same algebra, including its
multiplication. Put $\mathcal F=\Rep(U)_A$, the category of all
right $A$-modules. The algebra $A$ is Lagrangian in $\Rep(U)$,
and the canonical-algebra reconstruction therefore gives
\begin{equation}\label{eq:c-C6-defects}
 \begin{aligned}
 {}_{\one\boxtimes A}\Rep(C_6)_{\one\boxtimes A}
 &\simeq P_{39}^{\rm rev}\boxtimes{}_A\Rep(U)_A\\
 &\simeq P_{39}^{\rm rev}\boxtimes
                 \mathcal F\boxtimes\mathcal F^{\rm op}.
 \end{aligned}
\end{equation}
Here $\mathcal F^{\rm op}$ reverses tensor order. The total algebra
$\one\boxtimes A$ is not Lagrangian in $\Rep(C_6)$: its local
modules retain the pointed factor $P_{39}^{\rm rev}$.

For this original embedded extension, the result of
Chapter~\ref{ch:unitarity} is $\mathcal F\in\{\mathcal H_2,
\mathcal H_3\}$. No step here selects between them. In particular,
$A$ in this section continues to mean the original extension algebra,
not the separately selected canonical $\mathcal H_3$ algebra of
Chapter~\ref{ch:defects}. The transported multiplication preserves
whichever of these boundary categories the original embedding has.

The charge-six full state space can now be written explicitly.
Let $\mathcal H_{\lambda+N}$ be the Hilbert completion of the
lattice module $V_{\lambda+N}$, where
$\lambda\in N^*/N$. The ordinary diagonal lattice theory has
\begin{equation}\label{eq:c-C6-full-space}
 \HH_6=\bigoplus_{\lambda\in N^*/N}
     \mathcal H_{\lambda+N}\otimes
                     \overline{\mathcal H_{\lambda+N}},
 \qquad c_L=c_R=6.
\end{equation}
The bar denotes the conjugate antiholomorphic sector. There are
$|N^*/N|=\det N=39$ chiral sector labels, and its torus amplitude is
\[
 Z_6(\tau,\bar\tau)=\sum_{\lambda\in N^*/N}
       \left|\frac{\Theta_{\lambda+N}(\tau)}{\eta(\tau)^6}\right|^2.
\]
Positive definiteness of $N$ gives a unique zero-energy vector,
and the standard lattice construction supplies the positive local
OPE and sewing. This is an actual compact full theory, rather than
only a proposed central charge or branching function.

To interpret \eqref{eq:c-C6-defects} as physical defects on
\eqref{eq:c-C6-full-space}, one uses the positive rational extension
and bimodule-to-defect correspondence of
Chapters~\ref{ch:unitarity} and~\ref{ch:defects} for the actual
inclusion $C_6\subset V_N$. Its analytic premise is compatibility
of the unitary sector realization with this extension and its
intertwiners. Regularity of $C_6$ and the lattice bounded-insertion
theorem already give a finite-index rational subnet; that fact
alone should not be substituted for the required sector comparison.
With this analytic correspondence in place, the unit pointed
spectator in \eqref{eq:c-C6-defects} transports the original
$\mathcal F\boxtimes\mathcal F^{\rm op}$ defects, including their
junctions, to the charge-six theory. These defects preserve $C_6$.
The simple defects required to preserve the larger lattice algebra
$V_N$ pointwise remain pointed. This distinction also explains why
retaining the smaller preserved algebra is essential here.

The multiplication is the decisive extra information. Restricting the
matrix representing a line on a few states would not prove that its
junctions still associate correctly. Here the extension correspondence
transports the algebra that encodes those junctions.

\section{The reverse coset gives six vacua if labels are retained}

The same calculation explains why deleting the charge-six factor does
not automatically give the desired charge-two theory. The reverse
commutant in $E_8$ is exactly $V_{L_0}$. Write the original algebra
as
\[
 A=\one\oplus X\oplus2Y.
\]
For the $C_6$ labels indexed by a lattice charge $n$ and an extension
label $a$, the complementary branching functions are
$b_{n,a}=m_a\chi_n^{L_0}$ with multiplicities $m_a=1,1,2$.
An unreduced diagonal sum gives
\begin{equation}\label{eq:c-six-vacua}
 \sum_{n,a}|b_{n,a}|^2
   =(1^2+1^2+2^2)\sum_n|\chi_n^{L_0}|^2=6Z_{L_0}.
\end{equation}
Its vacuum coefficient is six. Identifying the repeated branching
spaces instead gives the ordinary unique-vacuum torus theory, but
then its rational category is the pointed $P_{39}$, not the
Haagerup center. The two tasks---removing repeated vacua and retaining
the noninvertible junction category---are independent.

An ordinary rational full theory preserving the lattice chiral algebra
$V_{L_0}$ has a weakly integral defect category: every simple defect
has integer squared Frobenius--Perron dimension. Here weak integrality
means that the global Frobenius--Perron dimension is an integer;
the assertion about its simple dimensions is a theorem for such
fusion categories. A unitary strong
tensor action of $\mathcal H_3$ would preserve the simplicity of
$\rho$ by Proposition~\ref{prop:c-faithful-descent} below. Its
dimension has irrational square, $d^2=(11+3\sqrt{13})/2$,
so this rational lattice interpretation cannot carry the desired
lattice-preserving Haagerup action. A new current-breaking action
would be a separate construction, not a consequence of
\eqref{eq:c-six-vacua}.

\section{What ordinary invariant cosets cannot remove}

There are several useful obstructions with carefully bounded scope.
The retained current Lie algebra in the accepted charge-eight
realization is $U_1=\mathfrak{sl}_2\oplus\mathfrak{sl}_2$, generating
the affine subalgebra $A_{1,1}\otimes A_{1,39}$. The currents of any
transparent compact denominator lie in this retained Lie algebra.
Positivity of its complementary
Sugawara stress gives
\[
 c_{\mathrm{denominator}}\leq\frac{158}{41},\qquad
 c_{\mathrm{coset}}\geq8-\frac{158}{41}=\frac{170}{41}>2.
\]
Allowing each original Haagerup line to transmit the denominator by
one group automorphism does not evade this bound. If a group label
$g_x$ is assigned to each simple line and every fusion summand $z$
of $x\otimes y$ obeys $g_z=g_xg_y$, then $\rho^2$ containing
both $1$ and $\rho$ forces
$g_\rho^2=1=g_\rho$. Thus $g_\rho=1$, and the other summands
force $g_\alpha=1$. All labels are trivial.

The statement assumes one automorphism per simple line and compatible
transmission through fusion junctions. General noninvertible maps that
mix or annihilate fields need not have this form. Generalized quotients
that produce new lines after gauging are also outside the argument.

The standard $A_2$ and $E_6$ stresses are not in $U$. A particularly
short diagnostic is to apply the $A_2$ stress operator $L_0^{A_2}$
to $H_R$: it selects the $A_2$ Cartan part of $h$, which is not in
the two-dimensional retained Cartan span. If the stress belonged to
$U$, closure would put that extra current in $U_1$, contradicting the
six-current theorem.

Finally $C_{6,1}=0$ and its exact weight-two algebra has dimension
seven. Its polynomial Virasoro-vector equations have no solutions of
charges one or two, even over $\C$. Complementation inside its
charge-six stress also excludes charges five and four. Thus a further
ordinary charge-four conformal coset of this particular $C_6$ cannot
produce charge two. These are exclusions of specified mechanisms,
not a classification of all possible Haagerup CFTs at charge two.

\section{A nearby rational stress from diagonal affine currents}

The compact diagonal of $A_{1,1}\otimes A_{1,39}$ has level forty.
This is the diagonal affine construction used in
Example~\ref{ex:voa-Ising-branching}: the current double poles add
the levels, and the coset stress subtracts the diagonal Sugawara
stress from the two numerator stresses.
In the coordinate gauge of Chapter~\ref{ch:lattice_embedding}, its
triple is
\[
 (H_L+H_R,\ 8640E_L+E_R,\ F_L/8640+F_R).
\]
The reciprocal rescaling makes this a PCT-stable compact diagonal.
Its GKO stress
\[
 e_0=\omega_L+\omega_R-\omega_{A_{1,40}}
\]
has charge
\[
 1+\frac{117}{41}-\frac{120}{42}
 =\frac{286}{287}=1-\frac6{41\cdot42}.
\]
It is the unitary minimal-model stress $M(41,42)$.
Add the diagonal Cartan stress, of charge one, to obtain the exact
near-two value
\[
 \frac{573}{287}=2-\frac1{287}.
\]
The vector $u_0=(\ell+h)/2=v-\ell$ has norm twenty and its
exponential has the minimum compatible affine energy ten. Its GKO
weight is therefore zero. The rational algebra
\[
 M(41,42)\otimes V_{\langle20\rangle}
\]
is actually embedded with this stress.

This is an existence statement for a rational chiral algebra near
charge two, not for a Haagerup full theory. Its rational categorical
dimension field is ramified only at $2,3,7,41$, while the Haagerup
dimension requires $\sqrt{13}$. Within the ordinary finite rational
framework preserving this algebra, that is an obstruction. Merely
preserving its stress would be a weaker premise.

\begin{example}[A discriminant detects the missing Haagerup dimension]
\label{ex:coset-ramification}
\index{ramification}\index{discriminant!number field}
Here is the arithmetic behind that obstruction. For a number field
$K$, the field trace $\operatorname{Tr}_{K/\Q}$ sums the images
under its complex embeddings. If $b_1,\ldots,b_r$ is an integral
basis of its ring of integers, its discriminant is
$\det(\operatorname{Tr}_{K/\Q}(b_i b_j))$.
A rational prime is \emph{ramified} precisely when it divides
this field discriminant. We use two standard facts: a ramified
prime of a subfield also ramifies in the larger field, and the
only primes that can ramify in $\Q(\zeta_N)$ divide $N$.

For $\Q(\sqrt{13})$, an integral basis is $1,\xi$ with
$\xi=(1+\sqrt{13})/2$ and $\xi^2=\xi+3$. Consequently
\[
 \operatorname{disc}\Q(\sqrt{13})
 =\det\begin{pmatrix}2&1\\1&7\end{pmatrix}=13.
\]
This is a discriminant of a number field, not the Gram determinant
of an $E_6$ or $A_2$ root lattice.

The minimal-model quantum dimensions above belong to
$K=\Q(\zeta_{82},\zeta_{84})=\Q(\zeta_{3444})$;
the norm-twenty lattice contributes only dimension-one simples.
Since $3444=2^2\cdot3\cdot7\cdot41$, the prime thirteen does
not ramify in $K$. Hence $K$ cannot contain $\Q(\sqrt{13})$,
or the number $d=1+\xi$.

For any simple separable Frobenius sewing algebra $A$ in the
stated representation category $\mathcal C$, with $A$ simple
as an $A$--$A$ bimodule, the dimension of a defect bimodule $X$ is
\[
 \FPdim_{{}_A\mathcal C_A}(X)
   =\frac{\FPdim_{\mathcal C}(X)}{\FPdim_{\mathcal C}(A)}\in K.
\]
Here numerator and denominator refer to underlying objects;
direct sums also have dimension in $K$. Thus neither a simple
nor a reducible defect preserving this rational algebra can
have Haagerup dimension $d$. No approximation to $c$ enters
the argument. The same arithmetic will filter the ordinary
$W_3$ family in Section~\ref{sec:identification-W3}.
\end{example}

\section{An exactly solvable four-dimensional stress algebra}

Now consider the commutant of the whole diagonal $A_{1,40}$ inside
$U$. Its weight-two space is four dimensional. Besides $e_0$ above,
it contains the two stresses $e_1,e_2$ of charges $2400/1271$ and
$70/31$, and a PCT-real diagonal singlet $x$ from the current
representation $(2,2)$. Normalize $\langle x,x\rangle=1$.

Use the Griess product $a*b=a_{(1)}b$. Put
\[
 f_i=\langle e_i,e_i\rangle,
 \qquad f=\left(\frac{143}{287},\frac{1200}{1271},\frac{35}{31}\right),
 \qquad\lambda=\left(\frac{11}{41},\frac{192}{1271},\frac{49}{31}\right).
\]
The complete multiplication table is
\begin{equation}\label{eq:c-axis-table}
 e_i*e_j=2\delta_{ij}e_i,
 \quad e_i*x=\lambda_i x,
 \quad x*x=\sum_i\frac{\lambda_i}{f_i}e_i.
\end{equation}
The last coefficients are forced by invariance of the form:
$\langle x*x,e_i\rangle=\langle x,x*e_i\rangle=\lambda_i$.
Representation parity excludes an $x$ component. Exact ambient OPE
identities verify that no omitted weight-two component remains.

For $e=\sum_i a_ie_i+tx$, the Virasoro idempotent equation $e*e=2e$
becomes
\begin{equation}\label{eq:c-axis-equations}
 a_i(1-a_i)=q_it^2,\qquad
 t(\lambda\cdot a-1)=0,
 \qquad q=\left(\frac7{26},\frac2{25},\frac7{10}\right).
\end{equation}
Its central charge is $c(e)=2f\cdot a$. A real solution has
$0\leq a_i\leq1$. These small equations describe actual stresses
because the quasiprimary and remaining Virasoro identities are
controlled in this current-free commutant.

At $c(e)=2$ and $t\ne0$, the linear equations give
$a_0=41/44-30a_1/11$ and $a_2=93/196+18a_1/49$.
Writing $b=a_1$ leaves two necessary quadratics,
\[
 2500b^2+620b-39=0,
 \qquad10676b^2-10820b+309=0.
\]
The first factors as $(250b-13)(10b+3)$. Its roots are $13/250$
and $-3/10$. Substitution in the second polynomial gives, respectively,
$-3512064/15625$ and $112896/25$, both nonzero. Thus no common
root exists, even over $\C$. The $t=0$ subset stresses also have
no charge two. This is a complete exclusion in this four-dimensional
slice, not in all of $U_2$.

The nonzero-$t$ solutions do give four real irrational charges:
\[
 1.537523906\ldots,\quad2.025004294\ldots,
 \quad3.117852848\ldots,\quad3.605333236\ldots.
\]
They occur in complementary pairs summing to $36/7$, the charge of
the diagonal-affine commutant. The one closest to two is
\begin{equation}\label{eq:c-cstar}
 c_*=\frac{18}{7}
 -4\sqrt{\frac{4166013889-141778560\sqrt{274}}{97483514569}}.
\end{equation}
\index{irrational stress!near charge two}
If $z=c_*-18/7$, then
\[
 4776692213881z^4-6532309777952z^2+1524572989696=0.
\]
Its displayed squared offset is irrational, so $c_*$ is not exactly
two. Rational $C_2$-cofinite VOAs have rational central charge;
therefore a chiral algebra with this stress cannot be strongly
rational, nor can a finite extension with the same stress become
strongly rational.

\section{What the irrational commutant actually contains}

Choose the stress $e$ of charge $c_*$ and put
\[
 f=\omega_{E_8}-e,\qquad
 N_e=\Com(\Vir_f,V_{E_8}),\qquad
 M_f=\Com(N_e,V_{E_8}).
\]
The vacuum-like-vector lemma proves
\[
 M_f=\Com(\Vir_e,V_{E_8}),\qquad
 N_e=\Com(M_f,V_{E_8}).
\]
Indeed a vector killed by the translation $e_{(0)}$ is killed by
all nonnegative modes of the whole $N_e$ algebra, not only the
Virasoro modes. This yields the first equality; the second follows
from the defining commutant and the reverse containment.

The diagonal affine algebra belongs to $M_f$. Write
\[
 f=\omega_{A_{1,40}}+g,
 \qquad c_g=\frac{36}{7}-c_*.
\]
The exact current decomposition under diagonal $SU(2)$ is
\[
 248=10(3)+10(5)+10(7)+6(9)+4(11).
\]
There is no singlet. Positivity of the complementary affine energy
therefore gives $(N_e)_1=0$.

\begin{constructioninput}[Vacuum kernels of the irrational factor]
Complete exact low-grade calculations of the ambient complementary
energy kernels give
\begin{equation}\label{eq:c-N-lowgrades}
 \bigl(\dim(N_e)_0,\ldots,\dim(N_e)_4\bigr)=(1,0,1,1,2).
\end{equation}
Thus the vacuum algebra agrees with its Virasoro vacuum module
through weight four. Its weight-four basis is
$L^e_{-4}\one,(L^e_{-2})^2\one$. Their Gram determinant is
$c_*^2(5c_*+22)/2>0$, so the two states are independent.
The opposite algebra has the exact finite kernel
\[
 (M_f)_2=(A_{1,40})_2\oplus\C g,
 \qquad\dim(M_f)_2=10.
\]
\end{constructioninput}
There is no additional spin-three or spin-four vacuum generator in
$N_e$. This rules out ordinary $W_3$ for this factor. It does not
prove $N_e=\Vir_e$: a first additional generator may occur later.
Similarly the opposite-algebra statement is not the all-grade equality
$M_f=A_{1,40}\otimes\Vir_g$.

\section{Three inequivalent proposals for the reduced state space}

The smallest attempted reduction projects onto the complementary
Virasoro vacuum on both chiralities. Its partition function would be
$|\chi_{N_e}|^2$. It cannot be modular invariant. If it were, the
ratio $\chi_{N_e}(-1/\tau)/\chi_{N_e}(\tau)$ would have constant
modulus one and hence be constant; at $\tau=i$ positivity makes the
constant one. The one-dimensional modular representation would have
$S=1$, $T=e^{-2\pi i c_*/24}$. The relation $(ST)^3=S^2$ would
force $c_*\in8\Z$. Additional sectors are therefore mandatory.

A second proposal branches only over $\Vir_f$,
\[
 V_{E_8}=\bigoplus_h B_h\otimes L(c_f,h),
 \qquad\mathcal H_{\rm naive}=\bigoplus_hB_h\otimes\overline{B_h}.
\]
This produces infinitely many reduced vacua. Already the three
diagonal affine currents are complementary Virasoro primaries of
weight one and have $e$-energy zero. Pairing their multiplicity space
with itself adds nine vacuum states. More generally $A_{1,40}$ has
infinitely many Virasoro primaries: its high-temperature growth has
effective charge $20/7$, whereas a finite sum of irreducible
Virasoro characters at that charge grows with effective charge at
most one. Every such affine primary still has zero $e$-energy.

The more careful third proposal branches over the full opposite
algebra,
\[
 V_{E_8}=\bigoplus_\lambda B_\lambda\otimes(M_f)_\lambda,
 \qquad\mathcal H_{\rm red}
      =\bigoplus_\lambda B_\lambda\otimes\overline{B_\lambda}.
\]
The unitary ordinary branching is an algebraic countable semisimple
\index{ordinary branching!without rationality}
sum, even though the two factors need not be rational. To see why,
the commuting product $N_e\otimes M_f$ has the full ambient stress.
Its representation on $V_{E_8}$ is therefore unitary with
finite-dimensional nonnegative integer ambient grades. On the
least remaining grade, split the finite-dimensional star-closed
zero-mode algebra into irreducible top spaces. A simple top
generates a simple unitary module: any invariant submodule has a
grade-by-grade invariant orthogonal complement; if the submodule
misses the simple top, that complement contains the top and hence
everything it generates. Removing these summands and repeating
grade by grade exhausts the representation.

The finite-dimensional simple tops for the commuting product
factor into a simple top for each factor. Their generated modules
are tensor products of ordinary simple modules. Each separate
graded piece is finite: tensoring it with one fixed nonzero
opposite ground vector embeds it in a single finite ambient grade.
There are only countably many such summands. Grouping identical
$M_f$ modules produces the multiplicity spaces $B_\lambda$ above;
fixing their opposite ground weight again proves finite-dimensional
grades and lower boundedness for each $B_\lambda$. All these
statements concern algebraic finite-energy vectors, so no direct
integral or completion has been silently inserted.

This prescription has exactly one zero-energy vector, because
all zero-$e$-energy states lie in the full opposite vacuum algebra.
But deleting the complementary energies may bring infinitely many
different labels into a bounded reduced-energy interval. Ambient
trace convergence does not prove reduced trace convergence.

It also does not permit selecting attractive branches at will.
Ambient currents have positive $e$-weights near
\[
 .0597642,\quad .0609687,\quad .0624674,\quad .0672163.
\]
Their complementary weights are between zero and one, so they are
tops for both factors. The full-opposite diagonal prescription
therefore retains scalar dimensions twice all these values, including
the first three below $.134433$. Ambient $SU(2)$ multiplicities are
not automatically reduced bulk multiplicities, but the existence of
these branches is unavoidable for this specified pairing.

\section{What sewing must still establish}
\index{sewing!coset completion}

An actual reduced local CFT needs all of the following mathematical
data: a complete branching with finite reduced degeneracies or a
specified continuous measure; a convergent torus amplitude and its
modular transformation; a positive local bulk OPE obeying sewing; and
a transport of Haagerup defects that preserves junction compositions
and associators. Integral spin from pairing matching complementary
labels is only one necessary check.

The next two worked arguments turn the convergence warning into
precise mathematical tests. Their additional hypotheses are part of
their statements; the low-grade calculations do not establish them.

\subsection{A conditional divergence theorem}

\begin{proposition}\label{prop:c-diagonal-divergence}
\index{thermal trace!conditional divergence}
Assume the all-grade identity $M_f=A_{1,40}\otimes\Vir_g$.
For the complete full-$M_f$ branching, the diagonal reduced state
space has divergent thermal trace at every positive inverse
temperature. The conclusion does not require $N_e=\Vir_e$ or
irreducibility of the multiplicity spaces $B_\lambda$.
\end{proposition}
\begin{proof}
Let $c_A=20/7$ and let $\beta,\epsilon>0$ be dimensionless inverse
temperatures conjugate to the two Virasoro energies. Thus their
modular parameters are $\tau_e=i\beta/(2\pi)$ and
$\tau_f=i\epsilon/(2\pi)$, with $q_e=e^{-\beta}$ and
$q_f=e^{-\epsilon}$. For a physical inverse temperature $b$ on a
circle of length $L$ and propagation speed $v$, the corresponding
dimensionless parameter is $2\pi vb/L$, in units $\hbar=k_B=1$.
Here the two parameters can be varied independently as trace variables.
Define the normalized two-temperature trace
\[
 F(\beta,\epsilon)=\Tr_{V_{E_8}}
 e^{-\beta(L^e_0-c_*/24)-\epsilon(L^f_0-c_f/24)}.
\]
The two nonnegative commuting energies sum to the ambient energy,
so convergence follows by comparison with an ordinary $E_8$ trace
at $\min(\beta,\epsilon)$. The same estimate on compact subsets
of two upper half-planes proves normal convergence of the
vacuum two-variable trace. The Krauel--Miyamoto multivariable
modular theorem applies to orthogonal conformal vectors summing
to the stress of a regular ambient VOA, with this analyticity;
its vacuum clause supplies the descendant analyticity needed in
the theorem. Neither reduced factor is required to be rational.
The ambient $E_8$ algebra is holomorphic and has $S=(1)$, so the
theorem gives
\[
 F(\beta,\epsilon)=F(4\pi^2/\beta,4\pi^2/\epsilon).
\]
This is simultaneous inversion of two variables; it does not assert
separate modular invariance in either variable. Retaining the positive
vacuum term on the right gives
\[
 F(\beta,\epsilon)\geq
 \exp\left(\frac{\pi^2c_*}{6\beta}
                        +\frac{\pi^2c_f}{6\epsilon}\right).
\]

Under the stated product hypothesis, an opposite simple module is
one of the forty-one affine integrable modules tensored with a
unitary Virasoro module of lowest weight $h\geq0$. A Verma spanning
bound therefore gives, uniformly in $\lambda$,
\[
 \chi_{(M_f)_\lambda}(\epsilon)\leq U(\epsilon):=
 \max_j\chi_{A_j}(\epsilon)\,
 e^{\epsilon c_g/24}\prod_{n\geq1}(1-e^{-n\epsilon})^{-1}.
\]
Finite affine modular transformation and eta inversion imply
\[
 \limsup_{\epsilon\downarrow0}\epsilon\log U(\epsilon)
                      \leq\frac{\pi^2(c_A+1)}6.
\]
The oscillator contribution is one. It is an upper bound on character
growth, not a replacement of the actual central charge $c_g$ by one.

Write
$\chi_{B_\lambda}(\beta)=\Tr_{B_\lambda}
e^{-\beta(L^e_0-c_*/24)}$, including the same chiral Casimir shift.
If $R(\beta)=\sum_\lambda\chi_{B_\lambda}(\beta)$ were finite,
positive branching would imply
$F(\beta,\epsilon)\leq U(\epsilon)R(\beta)$.
Multiply logarithms by $\epsilon$ and let $\epsilon$ decrease to
zero. The lower and upper bounds force $c_f\leq c_A+1$,
contradicting $c_f=c_A+c_g$ and $c_g>1$. Hence $R(\beta)$
diverges for every $\beta>0$.

One must still justify divergence of the diagonal square sum; a
divergent sum need not have divergent squares in general. Here the
positive state-by-state character expansion gives
\[
 \chi_{B_\lambda}(\beta)^2\geq\chi_{B_\lambda}(2\beta),
\]
because the square contains every pair of a basis vector with
itself. Summing yields a lower bound $R(2\beta)=\infty$ for
the diagonal thermal trace.
\end{proof}

The proof isolates a general obstruction: a proposed complementary
algebra whose simple characters all have a uniform growth rate below
the complementary central charge cannot be discarded in this
all-branches diagonal fashion. Additional opposite fields might
change the uniform bound, or a different left/right pairing might
change the construction. Neither possibility is settled by vacuum
dimensions through weight four.

\subsection{A continuous diagonal density can be solved exactly}

An alternative is a noncompact theory with a positive continuum of
weights. It is important to speak of a partition \emph{density}:
a continuum direct integral need not have a literal finite
Hilbert-space trace. A direct integral is the continuum analogue
of a direct sum of sectors, with a measure replacing the discrete
sum over labels. Integrating their characters gives a partition
density only after specifying that measure and its normalization.
Suppose this density itself, not merely a
regulator-dependent finite part, obeys exact modular inversion.

\begin{proposition}[Strict diagonal continuum]\label{prop:c-continuum}
\index{diagonal continuum!modular obstruction}
Fix $c>1$. Pair only identical positive-weight irreducible unitary
Virasoro modules, with a positive measure $\mu$ whose Laplace
integral is finite at every positive argument. Assume no vacuum atom
and exact simultaneous modular inversion at independently positive
left and right inverse temperatures. Then, up to a positive overall
constant,
\[
 d\mu(h)=\frac{\mathbf1_{h>(c-1)/24}}{
                     \sqrt{h-(c-1)/24}}\,dh.
\]
In particular its effective central charge is one.
\end{proposition}
\begin{proof}
Put $A=(c-1)/24$. A positive-weight unitary irreducible at $c>1$
has character
\[
 \chi_h(\beta)=\frac{e^{-\beta(h-A)}}{\eta(\beta)},
\qquad\eta(\beta)=e^{-\beta/24}
                               \prod_{n\geq1}(1-e^{-n\beta}).
\]
There are no positive real degenerate highest weights in this
unitary range. For $1<c<25$, write $c=1+6Q^2$ and
$Q=b+b^{-1}$ with $|b|=1$. The Kac values
$[Q^2-(rb+s/b)^2]/4$ are nonreal when $r\ne s$ and
nonpositive when $r=s$. For $c\geq25$, take $b>0$;
then $rb+s/b\geq Q$ for positive integers $r,s$, again giving
nonpositive values. Thus a positive-weight module has the Verma
character above. The vacuum is the exception, with its
$L_{-1}\one=0$ relation, and is treated separately below.
Define $G(s)=\int e^{-s(h-A)}d\mu(h)$. Strict diagonality gives
\[
 Z(\beta,\bar\beta)=
                 \frac{G(\beta+\bar\beta)}{
                                      \eta(\beta)\eta(\bar\beta)}.
\]
Here $\beta$ and $\bar\beta$ are independently positive,
dimensionless parameters in the same convention as above. Normal
convergence gives the analytic continuation needed when deriving this
identity from a physical modular amplitude.

Set $s=\beta+\bar\beta$, $p=\beta\bar\beta$, and
$t=4\pi^2s/p$. Eta inversion,
$\eta(\beta)=\sqrt{2\pi/\beta}\eta(4\pi^2/\beta)$,
turns modular invariance into
\[
 G(t)=\frac{\sqrt p}{2\pi}G(s)=\sqrt{\frac{s}{t}}G(s).
\]
For fixed $s$, positive temperatures realize every
$0<p\leq s^2/4$. Thus this equation holds whenever
$st\geq16\pi^2$. For arbitrary $s_1,s_2$, choose a common large
$u$ with $s_i u\geq16\pi^2$. It follows that
$\sqrt{s_1}G(s_1)=\sqrt uG(u)=\sqrt{s_2}G(s_2)$.
Hence $G(s)=C/\sqrt s$.

The gamma integral identifies the positive measure:
\[
 \int_A^\infty\frac{e^{-sh}}{\sqrt{h-A}}\,dh
                          =\sqrt\pi\frac{e^{-As}}{\sqrt s}.
\]
Uniqueness of positive Laplace transforms rules out any other measure.
For an elementary justification, weight it by $e^{-s_0h}$ to make
a finite measure and push it to $x=e^{-h}\in[0,1]$. The Laplace
values at $s_0+n$ fix all polynomial moments. Density of polynomials
in continuous functions on $[0,1]$ fixes that measure uniquely.
Undoing the positive weight fixes the original measure.

Its threshold is $h_{\min}=A$, so
$c_{\rm eff}=c-24h_{\min}=1$.
\end{proof}

\begin{example}[The noncompact boson as a partition density]
\label{ex:c-noncompact-boson}
\index{partition density}\index{noncompact boson}
The compact boson gives a familiar way to obtain an inverse-square-root
measure. Use the radius convention of Chapter~\ref{ch:cft},
\[
 p_L=\frac mR+\frac{nR}{2},\qquad
 p_R=\frac mR-\frac{nR}{2},\qquad m,n\in\Z.
\]
For independently positive $\beta,\bar\beta$, its partition function is
\[
 Z_R(\beta,\bar\beta)=
 \frac{1}{\eta(\beta)\eta(\bar\beta)}
 \sum_{m,n\in\Z}
 e^{-\beta p_L^2/2-\bar\beta p_R^2/2}.
\]
The eta factors count oscillators and include the two chiral Casimir
shifts. We use the same dimensionless inverse temperatures and eta
normalization as in Proposition~\ref{prop:c-continuum}.

Take $R\to\infty$ at fixed positive temperatures. To check suppression
of winding even when the two temperatures differ, put
$s=\beta+\bar\beta$ and $a=\beta-\bar\beta$. Completing the square gives
\[
 \frac\beta2p_L^2+\frac{\bar\beta}2p_R^2
 =\frac s2\left(\frac mR+\frac{anR}{2s}\right)^2
  +\frac{\beta\bar\beta}{2s}n^2R^2.
\]
The momentum Gaussian sum is $O(R)$ uniformly in its shift, as follows
by comparison with its Gaussian integral, or by Poisson summation.
The second term therefore suppresses the entire $n\ne0$ contribution
after division by $R$. For $n=0$, the spacing of $p=m/R$ is $1/R$,
so the remaining sum is a Riemann sum:
\begin{equation}
 \begin{split}
 \mathcal Z_\infty(\beta,\bar\beta)
 &:=\lim_{R\to\infty}\frac{Z_R(\beta,\bar\beta)}R\\
 &=\frac{\displaystyle\int_{\R}e^{-sp^2/2}\,dp}
         {\eta(\beta)\eta(\bar\beta)}
 =\frac{\sqrt{2\pi}}
         {\eta(\beta)\eta(\bar\beta)\sqrt{\beta+\bar\beta}}.
 \end{split}
 \label{eq:c-noncompact-density}
\end{equation}
Dividing instead by the target-space circumference $2\pi R$ changes
only the overall constant. Eta inversion verifies exact simultaneous
modular inversion of this density.

The surviving primary weights are diagonal:
$h=\bar h=p^2/2$. Both signs of $p$ give the same Virasoro weight,
and changing variables gives
\[
 \int_{\R} f(p^2/2)\,dp
 =\sqrt2\int_0^\infty f(h)h^{-1/2}\,dh.
\]
Thus our normalization has $d\mu(h)=\sqrt2\,h^{-1/2}dh$ for $h>0$.
The singularity at zero is integrable and carries no atom.
At $c=1$, the Heisenberg Fock character is
$e^{-\beta h}/\eta(\beta)$. It agrees with the irreducible Virasoro
character at generic momentum. The exceptional degenerate weights
$h=j^2/4$, $j\in\Z_{\ge0}$, form a countable set of measure zero
for this continuous measure. Their distinct Virasoro decompositions,
including the special vacuum relations at $p=0$, do not change the
density integral.

There is still no literal finite Hilbert-space trace in the noncompact
limit. The zero-mode Hilbert space is $L^2(\R)$, and the heat operator
has an infinite spatial-volume trace. Equation~
\eqref{eq:c-noncompact-density} is a volume-normalized trace, equal
to $2\pi$ times the trace per unit target-space circumference,
obtained from the compact regulator. The finite-radius
$m=n=0$ sector contributes $1/R$ times its oscillator factor and
disappears from this density. Its constant zero-mode wavefunction has no
normalizable limit on the real line. A generalized zero-momentum
state and the local identity operation remain meaningful, but there
is no normalizable vacuum atom in the spectral measure.

The continuum threshold is zero, so $c_{\rm eff}=1-24\cdot0=1$.
For example at equal temperatures, eta inversion gives
\[
 \mathcal Z_\infty(\beta,\beta)
 \sim\frac{\sqrt\beta}{2\sqrt\pi}
       e^{\pi^2/(3\beta)}\qquad(\beta\downarrow0).
\]
The polynomial prefactor records the continuum; the exponential has
the common chiral effective charge one.

For $c>1$, the proposition forces the same Gaussian measure shape
after the shift $h=A+p^2/2$, where $A=(c-1)/24$. This comparison
explains the threshold density and again gives
$c-24A=1$. The compact-boson limit itself still has $c=1$.
Constructing a theory at $c=c_*$ would additionally require actual
fields, OPEs, and compatible sewing and defect data. This Gaussian
calculation supplies none of those additional data for a Liouville
or Haagerup construction.
\end{example}

At $c=c_*$, effective charge two would instead require
$h_{\min}=(c_*-2)/24$, about $.00104185$. The diagonal modular
density forces the larger threshold $(c_*-1)/24$, about $.0427085$.
Thus this precise continuum repair does not do the proposed job.
It also does not identify a Liouville interaction theory: characters
and a measure supply no OPE coefficients or defect junctions.

Adding a positive vacuum atom cannot fix the strict diagonal ansatz.
The vacuum character has the extra numerator $1-e^{-\beta}$.
After stripping eta factors, fix $\bar\beta>0$ and send
$\beta$ to zero. The original numerator stays bounded by Laplace
finiteness. Modular inversion would therefore make the transformed
numerator tend to zero after its square-root prefactor. But the
transformed positive vacuum term grows at least as
$\exp(4\pi^2A/\beta)$, a contradiction. This reasoning requires
positivity; cancellations in a signed formal expression would not
support it. Reducible branching blocks and asymmetric Virasoro
pairings remain outside the strict diagonal hypothesis.

\section{What a genuine defect descent would guarantee}

The inherited extension action preserves the two commutant spaces
because it intertwines the stress modes. This is weaker than a
tensor action on a sewn reduced theory. Restricted normalized
operations may cease to be extremal, and their compositions need
not supply the required fusion junction spaces. There is nevertheless
a useful positive criterion once a genuine tensor functor exists.

\begin{proposition}[Haagerup tensor descent cannot collapse simples]
\label{prop:c-faithful-descent}
Let $F$ be a unitary strong tensor functor from $\HH_3$ into a
locally semisimple rigid category with simple unit and
finite-dimensional Hom spaces on the image. Then $F$ is fully
faithful.
\end{proposition}
\begin{proof}
Take the full category $\mathcal D$ of simple summands of the images
of the six simples. It is a fusion category: products and duals stay
inside the finite list of image summands. The induced functor to
$\mathcal D$ is dominant, meaning every simple target is a summand of
an image. Its right adjoint gives a connected
central algebra $A=R(\one)$ and
\[
 \operatorname{Dim}\mathcal D
       =\frac{\operatorname{Dim}\HH_3}{\FPdim A}.
\]
If $A$ is not the unit, its central lift contains the unit once and
at least one nonunit central simple. The smallest such dimension in
$Z(\HH_3)$ is $3d$, so
$\FPdim A\geq1+3d=d^2$. Consequently
\[
 \operatorname{Dim}\mathcal D\leq3+\frac3{d^2}
                                      <\frac{5+\sqrt5}{2}.
\]
Every nonpointed fusion category has global dimension at least the
last number. At rank two this is the Fibonacci minimum from
$X^2=1+nX$, $n\geq1$. At rank at least three, a noninvertible
has squared dimension at least two and two further simples contribute
at least two, giving global dimension at least four.
Thus $\mathcal D$ would be pointed. But $F(\rho)$ still has
Frobenius--Perron dimension $d$, irrational, whereas every object of
a pointed category has integer dimension. This is impossible.
Therefore $A=\one$; the dominant functor is an equivalence onto
$\mathcal D$, proving full faithfulness into the target.
\end{proof}

The premise is substantial. A BRST reduction, a branching construction,
or a generalized quotient would have to build this tensor functor,
including its unit and junction compatibility. Checking the Haagerup
fusion polynomial for a few compressed operators does not meet that
premise. The proposition says what a correct construction would
preserve, rather than declaring that the missing construction has
already been carried out.

\section{A current test that survives changes of pairing}

A unitary theory
with unique normalizable vacuum and two independent commuting
Hermitian holomorphic currents has a rank-two Heisenberg stress of
charge two. Subtracting it at charge $c_*$ would leave a unitary
Virasoro vacuum with $0<c_*-2<1/2$, which is impossible. An explicit
negative-norm witness at residual charge $k$ is
\[
 -108L_{-6}\one-264L_{-4}L_{-2}\one
             +93L_{-3}^2\one+64L_{-2}^3\one,
\]
whose squared norm is
$12k(2k-1)(128k+6943)<0$ on that interval.
This is a finite Virasoro calculation, so it can be checked without
invoking a classification of unitary minimal models. In the ordered
vacuum basis
$L_{-6}\one,L_{-4}L_{-2}\one,L_{-3}^2\one,L_{-2}^3\one$,
the Gram matrix is
\[
 k\begin{pmatrix}
 35/2&5&18&24\\
 5&(5k+16)/2&14&(9k+48)/2\\
 18&14&8k+36&30\\
 24&(9k+48)/2&30&3(k+8)(k+16)/4
 \end{pmatrix}.
\]
For example,
$L_4L_{-2}^3\one=(48+9k)L_{-2}\one$ gives the second-row,
fourth-column entry after applying $L_2$.
Multiplying the matrix by $(-108,-264,93,64)$ on both sides
gives $3072k^3+165096k^2-83316k$, which factors as the
negative expression above. Quotienting null states cannot remove
a vector of strictly negative norm in a unitary representation.

This is a conditional criterion for rejecting $c_*$. The lattice
chain's nearly spin-one states have not been proved to be two such
currents, so the criterion cannot yet be applied as an unconditional
identification or exclusion.

The exact construction at charge eight, the rational charge-six coset,
and the irrational chiral stress are three distinct achievements.
The physical $\rho$-projector chain remains a fourth problem: one
needs a controlled map from its lattice operators and sectors to the
full continuum theory. Neither a matching central charge nor a
matching pair of low dimensions supplies that map.

\section{Exercises}

\begin{exercise}\label{ex:c-coset-1}
Verify the norm, index, and determinant calculation in
\eqref{eq:c-rank6-lattice}. Explain why the squarefree determinant
proves saturation of $L_0$ in its plane.

\par\smallskip\noindent{\footnotesize\normalfont
\hyperref[sol:ex:c-coset-1]{Solution on page~\pageref*{sol:ex:c-coset-1}.}}\par
\end{exercise}
\begin{exercise}\label{ex:c-coset-2}
Derive the six-vacuum factor in \eqref{eq:c-six-vacua}. What additional
information would be needed to interpret a multiplicity-six vacuum
as a full theory with a nontrivial vacuum algebra?

\par\smallskip\noindent{\footnotesize\normalfont
\hyperref[sol:ex:c-coset-2]{Solution on page~\pageref*{sol:ex:c-coset-2}.}}\par
\end{exercise}
\begin{exercise}\label{ex:c-coset-3}
Compute $c(A_{1,40})$ and the GKO charge $286/287$. Check that adding
one Cartan stress gives $573/287$, and compare both exact distances
from two with $c_*-2$.

\par\smallskip\noindent{\footnotesize\normalfont
\hyperref[sol:ex:c-coset-3]{Solution on page~\pageref*{sol:ex:c-coset-3}.}}\par
\end{exercise}
\begin{exercise}\label{ex:c-coset-4}
Derive \eqref{eq:c-axis-equations} from the Griess multiplication
table. Explain why the central charge is $2f\cdot a$ rather than
an arbitrary quadratic expression once the idempotent equations hold.

\par\smallskip\noindent{\footnotesize\normalfont
\hyperref[sol:ex:c-coset-4]{Solution on page~\pageref*{sol:ex:c-coset-4}.}}\par
\end{exercise}
\begin{exercise}\label{ex:c-coset-5}
Compute the Virasoro Gram matrix on
$L_{-4}\one,L_{-2}^2\one$ and its determinant. Use it to justify
the two-dimensional lower bound at weight four in
\eqref{eq:c-N-lowgrades}.

\par\smallskip\noindent{\footnotesize\normalfont
\hyperref[sol:ex:c-coset-5]{Solution on page~\pageref*{sol:ex:c-coset-5}.}}\par
\end{exercise}
\begin{exercise}\label{ex:c-coset-6}
Show directly from $\rho^2=1+\rho+\alpha\rho+\alpha^2\rho$
that any fusion-compatible group-valued assignment to simple objects
is trivial. Explain which transmission assumption is needed to apply
this algebraic observation to a current coset.

\par\smallskip\noindent{\footnotesize\normalfont
\hyperref[sol:ex:c-coset-6]{Solution on page~\pageref*{sol:ex:c-coset-6}.}}\par
\end{exercise}
\begin{exercise}\label{ex:c-coset-7}
Explain why finiteness below every ambient-energy cutoff need not
imply finiteness below every reduced-energy cutoff after one positive
commuting energy is discarded. Illustrate the issue using projections of a
rank-two integer lattice onto a line of irrational slope.

\par\smallskip\noindent{\footnotesize\normalfont
\hyperref[sol:ex:c-coset-7]{Solution on page~\pageref*{sol:ex:c-coset-7}.}}\par
\end{exercise}
\begin{exercise}\label{ex:c-coset-8}
Give positive numbers $a_n$ for which $\sum_n a_n$ diverges while
$\sum_n a_n^2$ converges. Why does the stronger character inequality
$\chi_\lambda(\beta)^2\geq\chi_\lambda(2\beta)$ nevertheless complete
the proof of Proposition~\ref{prop:c-diagonal-divergence}?
State where integrality of state multiplicities is used.

\par\smallskip\noindent{\footnotesize\normalfont
\hyperref[sol:ex:c-coset-8]{Solution on page~\pageref*{sol:ex:c-coset-8}.}}\par
\end{exercise}
\begin{exercise}\label{ex:c-coset-9}
Starting from $\eta(\beta)=\sqrt{2\pi/\beta}\,
\eta(4\pi^2/\beta)$, rederive the functional equation for $G$ in
Proposition~\ref{prop:c-continuum}. Explain why checking only equal
temperatures would not justify the step that compares arbitrary
$s_1$ and $s_2$ through one common large $u$.

\par\smallskip\noindent{\footnotesize\normalfont
\hyperref[sol:ex:c-coset-9]{Solution on page~\pageref*{sol:ex:c-coset-9}.}}\par
\end{exercise}
\begin{exercise}\label{ex:c-coset-10}
Prove that a nonpointed fusion category of rank at least three has
global dimension at least four. For rank two, show that
$X^2=\one+nX$ forces the nonpointed minimum
$(5+\sqrt5)/2$. Use these bounds to explain the role of
$\FPdim A\geq d^2$ in the tensor-descent proposition.

\par\smallskip\noindent{\footnotesize\normalfont
\hyperref[sol:ex:c-coset-10]{Solution on page~\pageref*{sol:ex:c-coset-10}.}}\par
\end{exercise}
